\begin{statement}{C.2.UN.44}{Proved}{The fundamental group and the component-incidence constraint of every torus fibre}
Let $f:U\to\Delta$ be a proper holomorphic degeneration of connected
complex two-tori, with smooth total threefold. Write
\[
 f^{-1}(0)=\sum_{i=1}^r m_iS_i,\qquad
 S=\bigcup_iS_i,\qquad d=\gcd(m_1,\ldots,m_r).
                                                        \tag{C.FG.1}
\]
All fundamental groups below are based; changing the paths to the
chosen base point changes the indicated presentations by their actual
change-of-base-point isomorphisms. Choose a lift $t$ of the positive
generator of $\pi_1(\Delta^*)$, put $\Lambda=\pi_1(F)=\Z^4$, and
write $T$ for conjugation by $t$ on $\Lambda$. There are vectors
$v_i\in\Lambda$, obtained from positively oriented meridians of the
actual components, for which
\[
 \Gamma:=\pi_1(S)
  \simeq
  \bigl(\Lambda\rtimes_T\langle t\rangle\bigr)
       \big/\langle\!\langle x^{v_i}t^{m_i}:1\le i\le r
                    \rangle\!\rangle .                \tag{C.FG.2}
\]
Here $x^v$ denotes the element $v$ of the additive lattice written
multiplicatively. Its image $A$ in $\Gamma$ is abelian, and there is
an exact sequence
\[
       1\longrightarrow A\longrightarrow\Gamma
          \longrightarrow\Z/d\Z\longrightarrow1 .    \tag{C.FG.3}
\]
In particular the index is exactly $d$ and $A$ is generated by at most
four elements. There is an exact integral description as follows.
Set $n_i=m_i/d$, choose integers $c_i$ with $\sum_i c_in_i=1$, and
put $v_0=\sum_i c_iv_i$. Then
\[
 A\simeq\Lambda/R_d,\qquad
 R_d=(T^d-I)\Lambda+
       \left\langle T^jv_i-n_iv_0:
                         1\le i\le r,\ 0\le j<d\right\rangle_{\Z}.
                                                        \tag{C.FG.4}
\]
For $d=1$, the whole group is abelian and this specializes to
\[
 \pi_1(S)\simeq
   \Lambda\big/\left((T-I)\Lambda+
                   \langle v_i-m_iv_0\rangle_{\Z}\right).
                                                        \tag{C.FG.5}
\]
The normalized base change $s^d=f$ is an actual connected degree-$d$
\emph{\'etale} cover of the total space near $S$. Its reduced central
support $S_d$ is connected, its restriction $S_d\to S$ is the
connected cyclic covering belonging to the subgroup $A$, and
$\pi_1(S_d)=A$ under that covering injection.

For every embedded resolution of the central divisor by smooth
vertical blowups, let $\mathcal D$ be the dual complex of its simple
normal-crossing support. There is a surjection
\[
                  \pi_1(S)\longrightarrow\pi_1(\mathcal D).
                                                        \tag{C.FG.6}
\]
Its target has an abelian normal subgroup of index dividing $d$,
generated by at most four elements and of rank at most two.
Moreover the cup-product map
$\bigwedge^*H^1(\mathcal D,\Q)\to H^*(\mathcal D,\Q)$ is injective.
When the resolved support has
no triple points, its connected incidence graph is a tree or has
exactly one independent cycle. In particular two components of an
actual simple normal-crossing fibre with no triple points can meet
in at most two connected curves.

For an actual simple normal-crossing support with no triple points,
write $e_1$ for the number of its connected double curves. The zero
Euler criterion of UN39 takes the following joint form:
\[
 \begin{gathered}
 \#\{
      \text{independent incidence cycles}\}\le1,\\
 e(S)=0\ \Longleftrightarrow\
       \sum_{D\ {\rm double}}g(D)=e_1+\frac{j(T)}2 .
 \end{gathered}
                                                        \tag{C.FG.7}
\]
Here $j(T)$ is unchanged by replacing the action on $H^1(F)$ by
the corresponding action on $H_1(F)$ or its inverse. If all $m_i$
are equal, $j(T)=0$. If $1$ is not an eigenvalue of $T$, a
no-triple incidence graph must be a tree.
These are constraints and exact group formulas
for actual degenerations; they include arbitrary multiplicities
and do not assert the existence of a smoothing of abstract gluing data.
\end{statement}

\begin{proof}
\par\medskip\noindent\textbf{The actual small-tube retraction.}
Shrink the disc so that zero is its only critical value, and take
a compact closed tube $U_\epsilon=f^{-1}(|z|\le\epsilon)$ inside
the original open total space. Put $\varphi=|f|^2$ and choose a
smooth Riemannian metric near this tube. The local analytic gradient
inequality, in the form of
\cite[Theorem (LI), p.~775]{Kurdyka1998}, gives near each point of
$S$ an exponent smaller than one. Compactness, a finite coordinate
cover, and uniform equivalence of the coordinate Euclidean metrics
with the chosen metric give constants $c>0$ and $0<\theta<1$ with
\[
                 \|\nabla\varphi\|\ge c\varphi^\theta
                   \quad\hbox{on }U_\epsilon\setminus S.
                                                        \tag{C.FG.8}
\]
Indeed take the maximum of the finitely many exponents, then shrink
to $0\le\varphi<1$; properness permits a whole smaller tube to be
contained in the resulting neighbourhood of $S$. There is no
critical point of $\varphi$ off $S$, since $df$ is surjective there.

On that complement use the vector field
$V=-\nabla\varphi/\|\nabla\varphi\|^2$. Along its flow the value
of $\varphi$ decreases at speed one. A trajectory starting at level
$a$ therefore exists up to every time smaller than $a$: each
positive closed level band is compact and the vector field is
smooth there. Its remaining length below level $b$ is at most
\[
          \int_0^b\frac{du}{c u^\theta}
                 =\frac{b^{1-\theta}}{c(1-\theta)}.    \tag{C.FG.9}
\]
It consequently has a limit in $S$. Let $H(x,q)$ be the flow for
time $q\varphi(x)$ when $q<1$, and its limiting point when $q=1$;
put $H(x,q)=x$ on $S$. The estimate (C.FG.9) proves continuity
at points of $S$ and uniform convergence as $q\to1$. For $q<1$
away from $S$, ordinary smooth dependence of flows proves
continuity. Thus $H$ is an actual strong deformation retraction
of $U_\epsilon$ onto $S$. This also proves the needed retraction
for every later resolved or covered tube. No invariant-cycle
surjectivity is involved.

\par\medskip\noindent\textbf{The exact normal subgroup of meridians.}
Write $U^*=U_\epsilon\setminus S$. Proper horizontal transport
over the punctured disc gives a smooth torus bundle. Its homotopy
sequence, and $\pi_2(\Delta^*)=0$, identify
\[
       \pi_1(U^*)=\Lambda\rtimes_T\langle t\rangle,
       \qquad \pi_1(U^*)\longrightarrow\pi_1(\Delta^*)=\Z
       \text{ sends }x^v t^a\text{ to }a .             \tag{C.FG.10}
\]
The splitting exists because an extension onto the infinite cyclic
group splits after choosing a lift of its generator.

General position makes every loop in the smooth six-manifold
$U_\epsilon$ avoid the divisor $S$. A null-homotopy of a loop from
$U^*$ can be made transverse to the smooth part of $S$ and disjoint
from its singular set: the latter has real dimension at most two,
so a transverse real two-disc misses it. Removing little discs
around the finitely many transverse intersections expresses the
original loop as a product of conjugates of positive or negative
meridians. Conversely each such meridian bounds its normal disc
in $U_\epsilon$.

One meridian for each global irreducible component suffices. The
smooth locus of an irreducible component, with its intersections
with other components removed, is path connected. This follows
also by taking its connected normalization resolution and moving
paths off the proper complex analytic curves and points. A path
there moves positive normal circles through positive normal circles,
giving conjugate based meridians in $U^*$. Near a chosen smooth
point of $S_i$ the actual equation is $f=u z^{m_i}$ with $u$ a
unit. The meridian has winding $m_i$, hence is exactly
$x^{v_i}t^{m_i}$ in (C.FG.10), for a uniquely determined $v_i$
after the indicated choices. The retraction and this kernel
calculation prove (C.FG.2).

The image $A$ of $\Lambda$ is normal and abelian. Quotienting
(C.FG.2) by this image leaves precisely the presentation
$\langle t\mid t^{m_i}=1\text{ for all }i\rangle$.
This is $\Z/d\Z$, which proves the exact index in (C.FG.3).
For later comparison, abelianizing the same actual presentation
also gives
\[
 H_1(S,\Z)\simeq
 \frac{\bigl(\Lambda/(T-I)\Lambda\bigr)\oplus\Z}
      {\langle ([v_i],m_i):1\le i\le r\rangle_{\Z}}.
                                                        \tag{C.FG.11}
\]

\par\medskip\noindent\textbf{The connected normalized root cover.}
Work momentarily in the open complex neighbourhood of the closed
tube. Since all divisor multiplicities of $f$ are divisible by
$d$, unique factorization in the regular local analytic ring
gives $f=u h^d$ near every point, including singular points of the
support. The normalization of $s^d=u h^d$ is the explicit map
\[
       \{v^d=u\}\longrightarrow\{s^d=u h^d\},
                         \qquad s=vh .                \tag{C.FG.12}
\]
Its source is finite and \'etale over the smooth original chart:
$v$ is a unit and $d v^{d-1}$ never vanishes. It is normal and is
isomorphic to the base change over $h\ne0$, hence is exactly that
base change's normalization. These normalization maps glue.

Over $U^*$ this is the cyclic covering given by the winding number
modulo $d$ in (C.FG.10). That homomorphism is onto, so the cover
is connected. Its complement in the normalization is a divisor;
the full normalized cover is therefore connected as well.
The tube retraction $H$ lifts uniquely to this actual covering.
Its lift starts at the identity, fixes the inverse image of $S$
pointwise, and at time one has image in that inverse image. It is
therefore a strong deformation retraction onto $S_d$, proving
connectedness of $S_d$ and the asserted covering injection.
The covering subgroup in $\Gamma$ is the kernel of the winding
modulo $d$ map, which is exactly $A$ by (C.FG.3). In particular
this construction does not presume that the cover of an individual
$S_i$ is connected or that it has degree one.

There is also a useful integral calculation using the punctured
part of this same cover. Its group is
$\Lambda\rtimes_{T^d}\langle w\rangle$, where $w=t^d$.
The kernel of its map to $\pi_1(S_d)$ is normally generated by
all the lifted meridians
\[
                x^{T^jv_i}w^{n_i},
                  \qquad 1\le i\le r,\quad0\le j<d .  \tag{C.FG.13}
\]
To see completeness of this list, use the coset representatives
$1,t,\ldots,t^{d-1}$ for the winding cover in (C.FG.10).
Conjugating the original normal generators by these representatives,
and then by elements of the covering subgroup, gives exactly its
kernel. Thus no choice of the connected components over $S_i$
has omitted a meridian.

The relations with $j=0$ imply $w^{n_i}=x^{-v_i}$.
Multiplying their $c_i$th powers gives $w=x^{-v_0}$. Consequently
the covering group is abelian, $(T^d-I)\Lambda$ is killed, and
the remaining relations are precisely
$T^jv_i-n_iv_0=0$. Conversely the map $x^v\mapsto[v]$,
$w\mapsto-[v_0]$ from this presentation to $\Lambda/R_d$
respects every semidirect and meridian relation. It is inverse to
the map induced by $\Lambda$. This proves (C.FG.4). For $d=1$
the same argument proves (C.FG.5); in particular it first proves
$t\in A$ and only then imposes $T=I$ on that quotient.
The Smith-normal-form comparisons in
\path{code/A_meridian_groups.py} check the two integral
presentations and their torsion in explicit test matrices; the
presentation is proved here for every actual monodromy.

\par\medskip\noindent\textbf{Resolution and the actual dual-complex map.}
A smooth vertical point or curve blowup induces an isomorphism on
the fundamental group of the total tube. For a curve centre, cut
out a small normal disc bundle. Its boundary is an $S^3$-bundle
over the curve, and the exceptional neighbourhood retracts to the
$\PP^1$-bundle over that same curve. The boundary inclusion
induces the same isomorphism on fundamental groups as the old
disc bundle, since both fibre fundamental groups are zero. For
a point centre the boundary is $S^5$ and the exceptional fibre
is $\PP^2$, again both simply connected. Van Kampen with the
unchanged complement proves the assertion for the actual blowdown
map. Iteration, followed by the two tube retractions, gives
\[
              \pi_1(S^{\rm res})\xrightarrow{\ \simeq\ }\pi_1(S).
                                                        \tag{C.FG.14}
\]
The old multiplicities are retained on strict transforms. Every
new multiplicity is the order of the old fibre equation along its
centre and is divisible by $d$, so their gcd remains $d$.

For an SNC support $E=\bigcup E_i$, form the space $P$ from the
disjoint components $E_i$, one $D\times[0,1]$ for each connected
double curve, and one two-simplex for each triple point. The faces
are glued by the actual intersection inclusions. There are actual
continuous maps
\[
                   E\ \xleftarrow{\ p\ }\ P
                       \xrightarrow{\ q\ }\mathcal D,             
                                                        \tag{C.FG.15}
\]
where $p$ forgets the simplex coordinate and $q$ collapses every
connected stratum to its associated cell. Compatible triangulation
of the finite analytic family of components and intersections makes
all the inclusions closed subcomplex inclusions; see
\cite[\S3, Theorem~2]{Lojasiewicz1964}. The map $p$ is a homotopy
equivalence. Explicitly, for two subcomplexes the double mapping
cylinder projects to their union by the homotopy-pushout comparison
for a cofibration. Inductively add one component, applying that
comparison also to its intersection with the previous union.
The intersection's mapping cylinders and the triple-point
triangles are precisely the pieces already included in $P$.
This proves the comparison for the finite union and supplies a
homotopy inverse of $p$.

Van Kampen on these mapping cylinders gives a particularly precise
description of $q_*$. Start with the vertex fundamental groups and
edge paths, impose the relations induced by the two inclusions of
each $\pi_1(D)$, and then the triangle boundary relations.
Killing the vertex groups leaves exactly the presentation of
$\pi_1(\mathcal D)$. Thus
\[
 \pi_1(\mathcal D)\simeq
 \pi_1(E)\big/
       \langle\!\langle\operatorname{im}\pi_1(E_i):i
                              \rangle\!\rangle .       \tag{C.FG.16}
\]
The map $q\circ p^{-1}$, where $p^{-1}$ is this homotopy inverse,
and (C.FG.14) give (C.FG.6). In particular it is an actual
topological quotient map on fundamental groups, including when
two components have several connected double curves.

The image of $A$ is an abelian normal subgroup whose quotient is
a quotient of $\Z/d\Z$. If there are no triple points,
$\mathcal D$ is a finite connected graph. Let its free-group rank
be $b$ and let the index of that abelian subgroup be $a$.
The corresponding connected graph covering has $a$ times as many
vertices and edges, hence free fundamental group of rank
$1+a(b-1)$. An abelian free group has rank at most one. Thus
$b\le1$, proving the tree-or-one-cycle assertion and its two-vertex
consequence.

\par\medskip\noindent\textbf{The rank-two bound, with the actual finite-cover comparison.}
For any of these tubes, inclusion of a nearby fibre followed by
the constructed retraction is a continuous map
$\sigma:F\to S$. Restriction to its fundamental group has image
$A$ of finite index. Therefore
\[
 \sigma^*:H^1(S,\Q)\longrightarrow H^1(F,\Q)
           \quad\hbox{is injective}.                  \tag{C.FG.17}
\]
Indeed a homomorphism $\Gamma\to\Q$ vanishing on $A$ factors through
the finite group $\Gamma/A$ and is zero. This proves the stated
direction of the cohomology map. The map $\sigma^*$ in all degrees
is a ring homomorphism because it comes from that actual continuous
map. Its degree-one image is $T$-invariant, as is also immediate
from the semidirect presentation.

Apply this to the resolved tube $W$ with central support
$E=S^{\rm res}$. The map $q\circ p^{-1}:E\to\mathcal D$ of
(C.FG.15) induces an injection on $H^1$ because its fundamental
group map is onto. The composite
\[
 H^*(\mathcal D,\Q)\longrightarrow H^*(E,\Q)
               \xrightarrow{\ \sigma^*\ }H^*(F,\Q)
                  =\bigwedge^*H^1(F,\Q)               \tag{C.FG.18}
\]
is again a ring homomorphism and is injective in degree one.
The exterior algebra on any linear subspace of $H^1(F,\Q)$
injects into this exterior algebra. Consequently cup product
on $\mathcal D$ is injective on
$\bigwedge^*H^1(\mathcal D,\Q)$. Since $\mathcal D$ is a
two-dimensional cell complex, its third cohomology vanishes,
and three independent degree-one classes would give a
contradiction. Thus $b_1(\mathcal D)\le2$.

This also bounds the rank of its finite-index abelian subgroup,
rather than only the rank of its abelianization. Put
$\Pi=\pi_1(\mathcal D)$ and let $B$ be the image of $A$ in $\Pi$.
The index $a=[\Pi:B]$ divides $d$. Take the actual connected
degree-$a$ covering $\mathcal D'\to\mathcal D$ belonging to $B$,
and pull it back along the continuous map
$W\to E\to\mathcal D$ just constructed. The pullback is connected
because $\pi_1(W)\to\Pi$ is onto. Its covering subgroup is the
inverse image of $B$ in $\Gamma$. The quotient $\Pi/B$ is cyclic
of order $a$, generated by the image of $t$, so this inverse
image is precisely the kernel of the winding number modulo $a$.
The uniqueness of a based covering with a given subgroup
identifies this pullback, over $W$, with the actual normalized
root cover
\[
                  s^a=f\circ\pi:W_a\longrightarrow W. \tag{C.FG.19}
\]
It is \'{e}tale by (C.FG.12), since $a$ divides every resolved
component multiplicity. Its lifted tube retraction has connected
central support $E_a$. The lifted map
$E_a\to\mathcal D'$ is onto on fundamental groups: its image is
exactly $B$, by the inverse-image description of the covering
subgroup. Regarded over the $s$-disc, $W_a$ has connected nearby
fibre equal to the actual torus $F_{s^a}$. Hence (C.FG.17)--(C.FG.18)
apply to this map and give $b_1(\mathcal D')\le2$. Since
$\pi_1(\mathcal D')=B$ is finitely generated abelian, this is
exactly $\operatorname{rank}B\le2$, as claimed.

If a no-triple incidence graph had one cycle, its nonzero
$H^1$ would inject by (C.FG.18) into the $T$-invariant part of
$H^1(F,\Q)$. Thus $1$ must be an eigenvalue of $T$ in that case.

Finally, for the original support already SNC with no triple
points, UN39 has no resolution-centre terms and gives
$e(S)=j(T)-2\sum_D(g(D)-1)$, proving (C.FG.7).
If all multiplicities are $m$, the full trace formula is
$\det(I-T^k)=m\sum_i e(S_i^\circ)$ when $m\mid k$ and zero
otherwise. Taking $k$ also divisible by the semisimple order makes
$T^k$ unipotent, so this determinant is zero. Hence the common
sum is zero and all traces vanish; in particular $j(T)=0$.
\end{proof}


% C.2.G.1 from part_C2_global.tex

